% ============================================================================
\clearpage
\section*{v4.9.0-r05 Source-Only: The Millennium Matrix on the Doublet, a Third Axis near Rotating Black Holes, and a Photon from Matter and Tachyon}
\addcontentsline{toc}{section}{v4.9.0-r05 Source-Only: The Millennium Matrix on the Doublet, a Third Axis near Rotating Black Holes, and a Photon from Matter and Tachyon}
\label{sec:v490-r05-layer-transport-kerr-photon}
% ============================================================================

\paragraph{Purpose.}
Three questions of the author follow r04.
\begin{enumerate}[leftmargin=2em]
\item Can the Millennium matrix of the main body be the doublet matrix of r04?  The author also asks for a more
meaningful name.
\item Can amplitudes, the missing third axis of r04, enter near a black hole?
\item Can a photon arise when a tachyon meets a forming massive particle, and could such a photon run back and create
the tachyon at its origin?
\end{enumerate}
The computations are \texttt{llm/k41\_schichttransport\_kerr\_photon.py}.  No value is compared with data.  Earlier
layers are not rewritten.

\noindent\textbf{Source-only status.}
This v4.9.0-r05 layer is source-only.  It does not build a PDF, does not assign a new DOI,
and does not authorize a public release.  The v4.9.0 Version DOI
\href{https://doi.org/10.5281/zenodo.23534954}{10.5281/zenodo.23534954} and the Concept DOI
\href{https://doi.org/10.5281/zenodo.19210728}{10.5281/zenodo.19210728} are unchanged by it.

\subsection*{v4.9.0-r05 the Millennium matrix on the doublet}
\label{subsec:v490-r05-millennium}

The main body asks three things of the Millennium matrix \(\mathcal M\).  It transports one object through its layers
(\S\ref{sec:millennium-hc}).  It commutes with the involution and with the compensator, condition (T4) of
Definition~\ref{def:tensor-structural-kernel}.  And the compensator is its spectral projector
(Theorem~\ref{thm:hc-spectral-projector}).

\paragraph{Theoremlet v4.9.0-r05.1.}
On the doublet of r04, with \(J=\sigma_x\) and \(C_{\HC}=\tfrac12(1+\sigma_x)\) acting from the right:
\begin{enumerate}[leftmargin=2em]
\item[(a)] The real \(2\times2\) matrices satisfying (T4) are exactly \(M=\lambda_+C_{\HC}+\lambda_-A\).  Together with
spatial rotations acting from the left, these are all transports of the doublet compatible with (T4).
\item[(b)] For \(\lambda_+\neq\lambda_-\), \(C_{\HC}\) is the spectral projector of \(M\) onto the eigenvalue \(\lambda_+\),
as Theorem~\ref{thm:hc-spectral-projector} requires.
\item[(c)] Of the transports of Reading~B, (T4) holds for spatial rotations, for \(J_B\) itself and for the duality
rotation by \(\pi\), that is \(C\).  It fails for every other duality rotation, for \(T\) and for boosts (r02, r04).
\item[(d)] The invertible members with determinant \(1\) are \(e^{t\sigma_x}=e^{t}C_{\HC}+e^{-t}A\).  They preserve
\(\operatorname{Re}s=E^2-B^2\), the matter--tachyon coordinate, and weight the two ports of r03 by \(e^{\pm t}\).
\end{enumerate}

\paragraph{Proof sketch.}
\(M\sigma_x=\sigma_xM\) is a linear condition with a two-dimensional solution space spanned by \(1\) and \(\sigma_x\).
\(e^{t\sigma_x}{}^{\mathsf T}\sigma_ze^{t\sigma_x}=\sigma_z\) since \(\sigma_x\sigma_z=-\sigma_z\sigma_x\).  Checked
numerically.  \(\square\)

\paragraph{Consequence.}
On the doublet the Millennium matrix can only be a \emph{port-weighting matrix}.  It leaves the compensator port and the
partner port invariant and multiplies them by two factors.  It is not the Lorentz transport of Reading~B, since boosts
violate (T4).  The doublet therefore realises the compensator part of \(\mathcal H_{16}\), not the full \(\mathcal M_{16}\).
The four-term layer decomposition \(\mathcal M_{\rm tri}+\mathcal M_{\rm alg}+\mathcal M_{\rm cl}+\mathcal M_{\rm proj}\)
and the index \(16\) have no counterpart on a two-dimensional doublet.  With \(t=\eta_{\rm sym}\) the port weights would be
the Doppler factors \(e^{\pm\eta}\) of way~(C) in r01.  That is arithmetic, not an identification.

\paragraph{Naming (author's request, 2026-10-21).}
From this layer on, \(\mathcal M_{16}\) is called the \emph{layer-transport matrix}.  The name says what
\S\ref{sec:millennium-hc} defines: one object carried through its layers.  Its restriction to the doublet is called the
\emph{port-weighting matrix}.  Symbols, labels and the wording of earlier layers stay unchanged, so ``Millennium matrix''
in earlier text means the same object.  A renaming of the earlier text and of the PDF keywords would rewrite released
layers.  It is not done here.

\subsection*{v4.9.0-r05 a third axis, and rotating black holes}
\label{subsec:v490-r05-third-axis}

\paragraph{Theoremlet v4.9.0-r05.2.}
\begin{enumerate}[leftmargin=2em]
\item[(a)] At a point, \(E\) and \(B\) are real, so the state matrix of r04 has no \(\sigma_y\) part anywhere, near a black
hole or not.
\item[(b)] For time-harmonic fields with complex amplitudes \(\hat E,\hat B\), the matrix of time averages
\(\hat G=\hat F^\dagger\hat F\) is Hermitian, and its \(\sigma_y\) part is \(-2\operatorname{Im}(\hat E^*\cdot\hat B)\).  For a plane
wave this is \(\pm1\) (relative to \(\operatorname{tr}\hat G\)) for the two circular polarisations and \(0\) for linear
polarisation.  The third axis is helicity, the quantity generated by the duality rotation.  With it the disk of r04
becomes a sphere, with light on the helicity axis.
\item[(c)] The gravitational field has the same structure.  Its Weyl curvature is a complex symmetric \(3\times3\) matrix,
gravito-electric plus \(i\) gravito-magnetic.  For a Schwarzschild black hole the only component
\(\Psi_2=-M/r^3\) is real.  For a Kerr black hole \(\Psi_2=-M/(r-ia\cos\theta)^3\) is complex.  \(\operatorname{Im}\Psi_2^2\), the
Pontryagin part of the curvature, is positive in one hemisphere, zero on the equator and negative in the other.
\end{enumerate}

\paragraph{Proof sketch.}
(b): \(\operatorname{tr}(\hat G\sigma_y)=-2\operatorname{Im}\hat G_{12}\); for \(\hat B=\hat z\times\hat E\) with
\(\hat E=(1,\pm i,0)/\sqrt2\) one gets \(\mp1\) up to orientation.  (c): standard form of \(\Psi_2\) in the Kinnersley tetrad;
the sign of \(\operatorname{Im}\Psi_2^2\) follows the sign of \(\cos\theta\).  Checked numerically for \(M=1\), \(a=0.9\),
\(r=3\).  \(\square\)

\paragraph{Reading.}
Amplitudes do not enter because a black hole is near.  They enter with waves and quantum states.  What a black hole adds
depends on its rotation.  A static one is mirror-symmetric and singles out neither helicity.  A rotating one puts its own
curvature off the real axis, on opposite sides in its two hemispheres, which is the gravitational counterpart of the two
sides of the photon in r02.1.  Whether this couples to the helicity axis of (b), and whether Hawking radiation of a
rotating black hole would then carry a helicity imbalance between the hemispheres, is open.  A quantitative statement is a
prediction and requires a pre-registration note.  Theorem~\ref{thm:core} treats the static case only.

\subsection*{v4.9.0-r05 a photon from matter and tachyon}
\label{subsec:v490-r05-photon}

\paragraph{Author's question (2026-10-21).}
A tachyon meets a massive particle that is just forming.  Instead of reaching a massive state, the encounter yields a
photon.  The tachyon develops a pull through the impact, so the photon is drawn along the tachyon's former path,
collides again and again, reaches the tachyon's origin, vanishes there, and thereby creates the tachyon.

\paragraph{Theoremlet v4.9.0-r05.3.}
\begin{enumerate}[leftmargin=2em]
\item[(a)] Let \(z_M\) be of matter type and \(z_T\) of tachyon type with \(s(z_T)=-s(z_M)\).  Then
\(s(z_M+z_T)=2\,z_M\cdot z_T\).  The superposition is a null field exactly when \(z_M\cdot z_T=0\).  For a static electric
field \(E\) and a static magnetic field \(B\) of equal strength, the superposition is null exactly when \(E\perp B\).  Its
energy then flows at the speed of light along \(E\times B\), perpendicular to both.
\item[(b)] At an angle \(\vartheta\) between \(E\) and \(B\), the superposition sits at \(b=(0,\cos\vartheta)\) in the disk of
r04.  \(\vartheta=90^\circ\) is the photon at the centre, and \(\vartheta=0\) is the dyon of r03.2 at the edge.
\item[(c)] Let a tachyon leave the origin \(O\) and meet the particle at \(K\), a distance \(d\) away, at speed \(v>1\).  A
photon running back from \(K\) to \(O\) arrives \(d(1+1/v)\) after the departure, and this arrival lies in the future light
cone of the departure in every frame.  The order of departure and meeting is frame-dependent (for observer speeds above
\(1/v\) it reverses), but the photon leg is not.
\end{enumerate}

\paragraph{Proof sketch.}
(a) Expand \(s\).  For \(z_M=E\), \(z_T=iB\): \(z_M\cdot z_T=iE\cdot B\).  Null fields have \(|E\times B|=u\) (r03.3).
(b) \(s=2i|E|^2\cos\vartheta\) and \(u=|E|^2\).  (c) The separation of departure \((0,0)\) and photon arrival
\((d/v+d,0)\) is timelike and future-directed; the tachyon leg is spacelike.  Checked numerically for \(v=1/\bsym\).
\(\square\)

\paragraph{Reading.}
The first part of the picture has a precise form.  Matter and tachyon can combine to a photon, but only with crossed
orientation.  With parallel orientation they give the dyon of r03.  The photon does not move along the tachyon's path.
It moves perpendicular to both inputs.  The last part, a photon that arrives at the origin and thereby creates the
tachyon, cannot hold for a photon running forward in time: it always arrives after the tachyon has left.  What does hold
is the familiar reinterpretation of tachyon exchange.  Since the tachyon leg is spacelike, ``emitted at \(O\), absorbed at
\(K\)'' and ``emitted at \(K\), absorbed at \(O\)'' are the same exchange seen by different observers.  A closed
self-creation would need a photon leg running backward in time.  Without one, the repeated encounter of the picture is a
cycle of period \(d(1+1/v)\), not a self-creation.  If such a cycle returns on the other side of the photon, the tachyon
created in the next cycle is the charge conjugate of the first (r02.1).  Reading~B has no dynamics (r10), so the pull of
the picture cannot be derived or excluded here.

\subsection*{v4.9.0-r05 register entries}
\label{subsec:v490-r05-register}

\begin{center}
\small
\begin{tabular}{@{}>{\raggedright\arraybackslash}p{0.34\textwidth}>{\raggedright\arraybackslash}p{0.58\textwidth}@{}}
\toprule
statement & status after v4.9.0-r05\\
\midrule
Millennium matrix, (T4), Thm.~\ref{thm:hc-spectral-projector} & on the doublet only \(\lambda_+C_{\HC}+\lambda_-A\) (port weighting); boosts and general duality rotations violate (T4); full \(\mathcal M_{16}\) not identified (r05.1)\\
name & \(\mathcal M_{16}\): layer-transport matrix; doublet restriction: port-weighting matrix; earlier wording unchanged\\
third axis of r04 & helicity, for complex amplitudes; Kerr curvature off the real axis with opposite sides in the two hemispheres; coupling open (r05.2)\\
photon from matter \(+\) tachyon & null exactly for crossed orientation, flux perpendicular to both; return photon cannot create the tachyon (r05.3)\\
\bottomrule
\end{tabular}
\end{center}

\subsection*{v4.9.0-r05 ledger and guard update}
\label{subsec:v490-r05-ledger-guard}

\[
R_{\rm lt}^{v490\text{-}r05}=(0,0,1)
\]
for (doublet form of the layer-transport matrix decided, third axis and photon question answered structurally,
identification of the full \(\mathcal M_{16}\) and the Kerr coupling open).

The layer spends no guard:
\[
R_{\rm glob}^{(v490\text{-}r05,+)}=(0,1,1,1,1),
\qquad
\Gamma_{\rm guard}^{v490\text{-}r05}=(1,1,1,1,1),
\qquad
\beta_4=1/256\ \text{unspent}.
\]
It does not imply exact finite \(C5\) zero closure, exact all-readout terminality,
\(C6/T2\), a global photon-clock/RH theorem, or global all-readout closure.

